\documentclass[11pt]{article}

\usepackage[T1]{fontenc}
\usepackage{amsmath,amssymb,amsthm}
\usepackage[margin=1in]{geometry}
\usepackage[colorlinks=true,linkcolor=blue,citecolor=blue,urlcolor=blue]{hyperref}

\newtheorem{theorem}{Theorem}
\newtheorem{proposition}[theorem]{Proposition}
\newtheorem{lemma}[theorem]{Lemma}
\newtheorem{corollary}[theorem]{Corollary}
\newtheorem{conjecture}[theorem]{Conjecture}
\theoremstyle{remark}
\newtheorem{remark}[theorem]{Remark}

\newcommand{\R}{\mathbb R}
\newcommand{\Pj}{\mathbb P}
\newcommand{\ip}[2]{\left\langle #1,#2\right\rangle}
\newcommand{\norm}[1]{\left\|#1\right\|}
\newcommand{\supp}{\operatorname{supp}}
\newcommand{\weakto}{\rightharpoonup}
\newcommand{\HH}{\mathcal H}
\newcommand{\cF}{\mathcal F}
\setlength{\emergencystretch}{2em}

\title{The q4 Cross Defect Forces a Pressure-Blind Commutator Anomaly}
\author{John Lawson and Codex}
\date{25 July 2026}

\begin{document}
\maketitle

\begin{abstract}
This round resolves the first microlocal audit of the conditional q4
projected-Oseen entrance.  Its primal--adjoint cross density obeys an
exact local conservation law whose current contains both the
Navier--Stokes pressure and the projected adjoint pressure.  Pure Fourier
localisation cancels both pressures and forces a signed transport
commutator flux of size \(c_0\) through arbitrarily high cutoffs.

The nonzero terminal cross measure is supported on an
\(\HH^1\)-null slice, hence is not in \(H^{-1}\).  Its preterminal
\(H^{-1}\) norm therefore diverges.  More sharply, a localized
mollification theorem shows that finite
\[
 \int_t^{T^*}
 \norm{u(s)}_{L^{3,\infty}}\norm{\nabla a(s)}_2^2\,ds
\]
would force the zero weak adjoint trace to be attained strongly.
Consequently this weighted integral is infinite on every terminal
interval, although the unweighted adjoint dissipation is finite.
The temporal gain
\(\nabla a\in L_t^{22/9,2}L_x^2\) would close this q4 branch.

The complete cubic cross current, including both pressures, approaches
only the spatial Lorentz endpoint \(L^{21/16,7/5}\); the direct
\(L^2\)-capacity argument therefore remains out of reach.  No Clay
alternative is proved.
\end{abstract}

\section{Epistemic status and input}

Every theorem below is conditional on the repository's energy-efficient
exact q4 branch and on the previously constructed projected-Oseen
entrance.  These antecedents are not known for arbitrary Clay data.
External mathematical review is pending.

Let \(u\) be smooth on \([0,T)\), and let
\[
 a\in
 L^\infty(0,T;L^2_\sigma(\R^3))
 \cap
 L^2(0,T;\dot H^1_\sigma(\R^3))
\]
solve the frozen-drift adjoint.  The input properties are
\begin{equation}
 \label{eq:pair}
 \ip{u(t)}{a(t)}=c_0>0,
 \qquad
 a(t)\weakto0
 \quad(t\uparrow T),
\end{equation}
and
\begin{equation}
 \label{eq:cross-measure}
 u(s_j)\cdot a(s_j)\,dx
 \overset{*}{\rightharpoonup}\zeta,
 \qquad
 \zeta(\R^3)=c_0,
 \qquad
 \supp\zeta\subset\sigma,
\end{equation}
where
\begin{equation}
 \label{eq:slice}
 \frac35\leq\dim_{\rm H}\sigma\leq1,
 \qquad
 \HH^1(\sigma)=0.
\end{equation}
The measures in \eqref{eq:cross-measure} have uniformly bounded total
variation and are tight.

The Clay problem remains unsolved.

\section{Exact local cross current}

The preterminal equations in pressure form are
\begin{align}
 \label{eq:ns}
 \partial_tu-\nu\Delta u+(u\cdot\nabla)u+\nabla p&=0,\\
 \label{eq:adj}
 \partial_ta+\nu\Delta a+(u\cdot\nabla)a+\nabla q&=0,
\end{align}
with both fields divergence free.  The whole-space pressures satisfy
\begin{equation}
 \label{eq:pressures}
 -\Delta p=\partial_i\partial_\ell(u_i u_\ell),
 \qquad
 -\Delta q=\partial_i\partial_\ell(u_\ell a_i).
\end{equation}

\begin{proposition}[Local primal--adjoint conservation]
\label{prop:local}
Put \(c:=u\cdot a\).  Then
\begin{equation}
 \label{eq:local-law}
 \partial_tc+\nabla\cdot J=0,
\end{equation}
where
\begin{equation}
 \label{eq:current}
 \boxed{
 J_k
 =
 u_kc+p\,a_k+q\,u_k
 -\nu\bigl(\partial_ku_i\,a_i-u_i\partial_ka_i\bigr).
 }
\end{equation}
For a scalar cutoff \(\phi\), the exact pressure payment on
\([t,s]\) is
\begin{equation}
 \label{eq:pressure-payment}
 \mathcal P_\phi[t,s]
 :=
 \int_t^s\int_{\R^3}
 (p\,a+q\,u)\cdot\nabla\phi\,dx\,d\tau.
\end{equation}
\end{proposition}

\begin{proof}
Differentiate \(c=u_i a_i\) and insert
\eqref{eq:ns}--\eqref{eq:adj}:
\[
 \partial_tc
 =
 \nu(\Delta u_i\,a_i-u_i\Delta a_i)
 -(u\cdot\nabla)c
 -\partial_i p\,a_i-u_i\partial_iq.
\]
The diffusion difference is
\[
 \partial_k(\partial_ku_i\,a_i-u_i\partial_ka_i).
\]
The remaining three terms are respectively the divergences of
\(-uc\), \(-pa\), and \(-qu\).  This proves
\eqref{eq:local-law}--\eqref{eq:current}.

Testing with \(\phi\) gives
\begin{equation}
 \label{eq:cutoff-law}
 \int\phi c(s)-\int\phi c(t)
 =
 \int_t^s\int J\cdot\nabla\phi,
\end{equation}
whose pressure component is \eqref{eq:pressure-payment}.
\end{proof}

\begin{remark}
The functional \(\mathcal P_\phi\) is signed.  Neither
\eqref{eq:local-law} nor the pressure Poisson equations give it a
coercive lower bound.
\end{remark}

Passing \(s\uparrow T\) in \eqref{eq:cutoff-law} yields
\begin{equation}
 \label{eq:terminal-current}
 \boxed{
 \zeta=c(t)\,dx-\nabla\cdot G_t,
 \qquad
 G_t:=\int_t^T J(s)\,ds,
 }
\end{equation}
in the distributional spaces established below.

\section{The \texorpdfstring{\(\HH^1\)}{H1}-null slice forces an
\texorpdfstring{\(H^{-1}\)}{H-1} cascade}

\begin{lemma}[Capacity endpoint]
\label{lem:capacity}
If \(S\subset\R^3\) satisfies \(\HH^1(S)=0\), every
\(H^{-1}(\R^3)\) distribution supported on \(S\) is zero.
\end{lemma}

\begin{proof}
Let \(F\in H^{-1}\) be supported on \(S\), and fix a test function
\(\psi\).  Zero \(\HH^1\) measure implies zero \(W^{1,2}\)-capacity in
three dimensions.  The compact set
\(\supp F\cap\supp\psi\subset S\) therefore has smooth neighbourhood
cutoffs \(\chi_n\) such that
\[
 \chi_n=1\ \hbox{near }\supp F\cap\supp\psi,
 \qquad
 \norm{\chi_n}_{H^1}\longrightarrow0.
\]
\[
 \ip{F}{\psi}=\ip{F}{\chi_n\psi}\longrightarrow0.
\]
Thus \(F=0\).
\end{proof}

\begin{theorem}[Forced microlocal divergence]
\label{thm:hminus}
The terminal cross measure satisfies
\begin{equation}
 \label{eq:not-hminus}
 \boxed{\zeta\notin H^{-1}(\R^3).}
\end{equation}
For every inhomogeneous Littlewood--Paley decomposition,
\begin{equation}
 \label{eq:dyadic-divergence}
 \sum_{k\geq-1}2^{-2k}\norm{\Delta_k\zeta}_2^2=\infty.
\end{equation}
Along the selected record sequence,
\begin{equation}
 \label{eq:preterminal-divergence}
 \boxed{
 \norm{u(s_j)\cdot a(s_j)}_{H^{-1}}
 \longrightarrow\infty.
 }
\end{equation}
\end{theorem}

\begin{proof}
Equations \eqref{eq:cross-measure}--\eqref{eq:slice} and
Lemma~\ref{lem:capacity} prove \eqref{eq:not-hminus}.
The Littlewood--Paley characterisation of \(H^{-1}\) gives
\eqref{eq:dyadic-divergence}.

Let
\[
 M(t):=\norm{u(t)}_{L^{3,\infty}},
 \qquad
 U_2:=\sup_{0<t<T}\bigl(\norm{u(t)}_2+\norm{a(t)}_2\bigr).
\]
Lorentz Sobolev and Lorentz H\"older imply
\begin{equation}
 \label{eq:hminus-upper}
 \norm{u(t)\cdot a(t)}_{H^{-1}}
 \lesssim U_2M(t).
\end{equation}
If \eqref{eq:preterminal-divergence} failed, a bounded subsequence
would converge weakly in \(H^{-1}\).  Its distributional limit is
\(\zeta\), contradicting \eqref{eq:not-hminus}.

For a finite dyadic range, tight weak-star convergence of the measures
also gives
\[
 \Delta_k(u(s_j)\cdot a(s_j))
 \longrightarrow\Delta_k\zeta
 \quad\hbox{strongly in }L^2.
\]
Finite partial sums then recover the same conclusion from
\eqref{eq:dyadic-divergence}.
\end{proof}

\begin{corollary}[A direct capacity-closing condition]
\label{cor:l2-current}
If the total current \(G_t\) in \eqref{eq:terminal-current} belongs to
\(L^2(\R^3)\) for one \(t<T\), then the q4 entrance is impossible.
\end{corollary}

\begin{proof}
Equation \eqref{eq:hminus-upper} gives \(c(t)\in H^{-1}\).
If \(G_t\in L^2\), then \(\nabla\cdot G_t\in H^{-1}\), so
\eqref{eq:terminal-current} would put \(\zeta\) in \(H^{-1}\).
This contradicts Theorem~\ref{thm:hminus}.
\end{proof}

\section{Pressure-free Fourier cascade}

Let \(B_N=m(D/N)\), where \(m\) is real, even, smooth, compactly
supported, and equal to one near the origin.  It is self-adjoint and
commutes with \(\Delta\) and \(\Pj\).

\begin{theorem}[Exact spectral commutator]
\label{thm:commutator}
For \(T_u:=u\cdot\nabla\),
\begin{equation}
 \label{eq:spectral-identity}
 \boxed{
 \frac d{dt}\ip{u}{B_Na}
 =
 \ip{[T_u,B_N]u}{a}.
 }
\end{equation}
For every \(t_0<T\),
\begin{equation}
 \label{eq:spectral-cascade}
 \boxed{
 \lim_{N\to\infty}
 \int_{t_0}^{T}\ip{[T_u,B_N]u}{a}\,dt
 =
 -c_0.
 }
\end{equation}
\end{theorem}

\begin{proof}
Differentiate \(\ip{u}{B_Na}\).  The diffusion terms cancel by
self-adjointness and commutation.  The Leray projections disappear
against solenoidal fields and commute with \(B_N\).  Since \(T_u\) is
skew-adjoint,
\[
 \frac d{dt}\ip{u}{B_Na}
 =
 -\ip{B_NT_uu}{a}-\ip{B_Nu}{T_ua}
 =
 \ip{T_uB_Nu-B_NT_uu}{a}.
\]
This is \eqref{eq:spectral-identity}.

For fixed \(N\), spatial tightness and weak convergence of \(u(s_j)\)
give
\[
 B_Nu(s_j)\longrightarrow B_Nu_*
 \quad\hbox{strongly in }L^2.
\]
Together with \(a(s_j)\weakto0\),
\[
 \ip{u(s_j)}{B_Na(s_j)}
 =
 \ip{B_Nu(s_j)}{a(s_j)}
 \longrightarrow0.
\]
Integrating \eqref{eq:spectral-identity} to \(s_j\) and taking
\(j\to\infty\) yields
\[
 \int_{t_0}^{T}\ip{[T_u,B_N]u}{a}\,dt
 =
 -\ip{u(t_0)}{B_Na(t_0)}.
\]
The fixed-\(N\) commutator is time integrable because its smooth kernel
and derivative map the \(L^1\) tensor \(u\otimes u\) into \(L^2\).
Finally \(B_N\to I\) strongly on \(L^2\), and
\eqref{eq:pair} proves \eqref{eq:spectral-cascade}.
\end{proof}

\begin{remark}
Both pressures cancel algebraically in
\eqref{eq:spectral-identity}.  The terminal defect is therefore a
transport-commutator anomaly in global frequency, even though pressure
remains present in every local spatial current.
\end{remark}

\section{Weighted renormalisation}

Reverse time by setting
\[
 \tau:=T-t,\qquad
 v(\tau):=a(T-\tau),\qquad
 b(\tau):=-u(T-\tau),\qquad
 Q(\tau):=-q(T-\tau).
\]
Then
\begin{equation}
 \label{eq:reverse}
 \partial_\tau v-\nu\Delta v
 +(b\cdot\nabla)v+\nabla Q=0,
 \qquad
 \nabla\cdot b=\nabla\cdot v=0,
\end{equation}
and \(v(\tau)\weakto0\) as \(\tau\downarrow0\).

\begin{theorem}[Weighted renormalisation criterion]
\label{thm:renormalisation}
Suppose \(b,v\in L^\infty_\tau L^2_x\cap
L^2_\tau\dot H^1_x\) solve \eqref{eq:reverse} on
\((0,\tau_0)\), \(v(\tau)\weakto0\), and
\begin{equation}
 \label{eq:weighted-finite}
 \int_0^{\tau_0}
 \norm{b(\tau)}_{L^{3,\infty}}
 \norm{\nabla v(\tau)}_2^2\,d\tau<\infty.
\end{equation}
Then \(v=0\) on \((0,\tau_0)\).
\end{theorem}

\begin{proof}
For a spatial mollifier \(\rho_\varepsilon\), define
\[
 R_\varepsilon
 :=
 (b\otimes v)_\varepsilon-b\otimes v_\varepsilon.
\]
For almost every \(\tau\),
\begin{equation}
 \label{eq:comm-bound}
 \norm{R_\varepsilon}_2
 \lesssim
 \norm{b}_{L^{3,\infty}}\norm{\nabla v}_2,
 \qquad
 R_\varepsilon\longrightarrow0
 \quad\hbox{in }L^2.
\end{equation}
The first estimate uses
\(\dot H^1\hookrightarrow L^{6,2}\); the convergence follows because
approximate identities are strong on \(L^2\) and \(L^{6,2}\).
Since
\(\norm{\nabla v_\varepsilon}_2\leq\norm{\nabla v}_2\),
\[
 \left|\int R_\varepsilon:\nabla v_\varepsilon\right|
 \lesssim
 \norm{b}_{L^{3,\infty}}\norm{\nabla v}_2^2.
\]
Dominated convergence and \eqref{eq:weighted-finite} remove the
interior transport commutator.

The weak initial trace must still be handled without assuming strong
continuity.  Fix a compact spatial cutoff \(\chi_R\).  For fixed
\(R,\varepsilon\), the map
\[
 f\longmapsto\chi_R(\rho_\varepsilon*f)
\]
is compact on \(L^2\).  Hence
\[
 \norm{\chi_Rv_\varepsilon(\tau)}_2\longrightarrow0
 \quad(\tau\downarrow0).
\]
The mollified equation is
\[
 \partial_\tau v_\varepsilon-\nu\Delta v_\varepsilon
 +\nabla\cdot(b\otimes v_\varepsilon)+\nabla Q_\varepsilon
 =-\nabla\cdot R_\varepsilon.
\]
Testing it on \((\delta,\tau_1)\) with
\(\chi_R^2v_\varepsilon\) gives
\begin{align*}
&\frac12\int\chi_R^2|v_\varepsilon(\tau_1)|^2
 +\nu\int_\delta^{\tau_1}\int
       \chi_R^2|\nabla v_\varepsilon|^2\\
={}&\frac12\int\chi_R^2|v_\varepsilon(\delta)|^2\\
&+\int_\delta^{\tau_1}\int
 \left[
 \frac{\nu}{2}|v_\varepsilon|^2\Delta\chi_R^2
 +\frac12|v_\varepsilon|^2b\cdot\nabla\chi_R^2
 +Q_\varepsilon v_\varepsilon\cdot\nabla\chi_R^2
 \right]\\
&+\int_\delta^{\tau_1}\int
 \left[
 \chi_R^2R_\varepsilon:\nabla v_\varepsilon
 +R_\varepsilon:(\nabla\chi_R^2\otimes v_\varepsilon)
 \right].
\end{align*}
First let \(\delta\downarrow0\) using the compact local weak trace, then
let \(\varepsilon\downarrow0\).  Equation \eqref{eq:comm-bound} and
dominated convergence remove the interior commutator.
When the derivative lands on \(\chi_R\), the commutator is bounded by
\[
 C_R\norm{b}_{L^{3,\infty}}\norm{\nabla v}_2\norm{v}_2.
\]
Energy interpolation puts
\(\norm{b}_{L^{3,\infty}}\) in \(L^4_\tau\), so this is time
integrable; strong convergence of \(R_\varepsilon\) makes the cutoff
commutator vanish as well.

No energy enters from spatial infinity.  Indeed, energy interpolation
gives \(b,v\in L^4_\tau L^3_x\), while
\[
 -\Delta Q=\partial_i\partial_\ell(b_\ell v_i)
\]
and Calder\'on--Zygmund theory give
\[
 Q\in L^2_\tau L^{3/2}_x.
\]
Thus
\[
 b|v|^2,\ Qv\in L^{4/3}_\tau L^1_x,
 \qquad
 |v|^2\in L^\infty_\tau L^1_x.
\]
All cutoff fluxes vanish as \(R\to\infty\).  The limiting identity is
\[
 \frac12\norm{v(\tau_1)}_2^2
 +\nu\int_0^{\tau_1}\norm{\nabla v(\tau)}_2^2\,d\tau
 =0.
\]
It holds for almost every \(\tau_1\leq\tau_0\).  Therefore \(v=0\)
as an energy-class solution.
\end{proof}

\begin{corollary}[Forced weighted adjoint dissipation]
\label{cor:weighted-divergence}
The q4 entrance satisfies
\begin{equation}
 \label{eq:weighted-divergence}
 \boxed{
 \int_t^T
 \norm{u(s)}_{L^{3,\infty}}
 \norm{\nabla a(s)}_2^2\,ds
 =\infty
 \quad\text{for every }t<T.
 }
\end{equation}
\end{corollary}

\begin{proof}
Finiteness on one terminal interval would satisfy
Theorem~\ref{thm:renormalisation} after reversing time, making
\(a=0\) there.  This contradicts \eqref{eq:pair}.
\end{proof}

\begin{corollary}[One sufficient temporal gain]
\label{cor:gain}
If, on one terminal interval,
\begin{equation}
 \label{eq:gain}
 \nabla a\in L_t^{22/9,2}L_x^2,
\end{equation}
then the q4 entrance is impossible.
\end{corollary}

\begin{proof}
The q4 clock gives
\[
 M(t):=\norm{u(t)}_{L^{3,\infty}}
 \in L_t^{11/2,\infty}.
\]
Equation \eqref{eq:gain} gives
\[
 \norm{\nabla a}_2^2\in L_t^{11/9,1}.
\]
Endpoint Lorentz H\"older makes the integral in
\eqref{eq:weighted-divergence} finite, a contradiction.
\end{proof}

\section{Optimised pressure-current hull}

Write
\[
 J=J^{\rm cub}+J^\nu,
\]
\[
 J^{\rm cub}:=u(u\cdot a)+pa+qu,
 \qquad
 J^\nu_k:=-\nu(\partial_ku_i\,a_i-u_i\partial_ka_i).
\]
Set
\[
 X(t):=\norm{\nabla u(t)}_2,\qquad
 Y(t):=\norm{\nabla a(t)}_2.
\]

\begin{proposition}[Exact interpolation hull]
\label{prop:hull}
For \(0\leq\theta<3/7\), define
\[
 \frac1{r_\theta}:=\frac13-\frac{\theta}{12},
 \qquad
 p_\theta:=\frac6{5-\theta},
 \qquad
 \rho_\theta:=\frac2{1+\theta}.
\]
Then
\begin{equation}
 \label{eq:cubic-hull}
 \boxed{
 \norm{J^{\rm cub}(t)}_{L^{p_\theta,\rho_\theta}}
 \lesssim
 M(t)^{2-\theta}X(t)^\theta Y(t),
 }
\end{equation}
and the right side is time integrable.  Also
\begin{equation}
 \label{eq:viscous-hull}
 \boxed{
 \norm{J^\nu(t)}_{L^{3/2,1}}
 \lesssim\nu X(t)Y(t),
 }
\end{equation}
with a time-integrable right side.
\end{proposition}

\begin{proof}
Real interpolation and Lorentz Sobolev give
\[
 \norm{u}_{L^{r_\theta,4/\theta}}
 \lesssim
 M^{1-\theta/2}X^{\theta/2},
 \qquad
 \norm{a}_{L^{6,2}}\lesssim Y,
\]
where \(4/\theta=\infty\) at \(\theta=0\).
Lorentz H\"older applied to \(u(u\cdot a)\) proves
\eqref{eq:cubic-hull} for the advective term.  The pressure equations
\eqref{eq:pressures}, Lorentz Calder\'on--Zygmund bounds, and the same
three factors prove it for \(pa\) and \(qu\).

The temporal reciprocal sum is
\[
 \frac{2(2-\theta)}{11}
 +\frac{\theta}{2}
 +\frac12
 =
 \frac{19+7\theta}{22}<1.
\]
Therefore the coefficient is in \(L^1_t\).  Equation
\eqref{eq:viscous-hull} follows from Lorentz Sobolev and
Cauchy--Schwarz in time.
\end{proof}

As \(\theta\uparrow3/7\),
\begin{equation}
 \label{eq:hull-endpoint}
 p_\theta\uparrow\frac{21}{16},
 \qquad
 \rho_\theta\downarrow\frac75.
\end{equation}
The dual spatial exponent required of a cutoff gradient approaches
\[
 \left(\frac{21}{16}\right)'=\frac{21}{5},
\]
while the zero-\(\HH^1\) capacity cutoffs provide small \(L^2\)
gradients.  The viscous term still asks for the \(L^{3,\infty}\)
gradient endpoint.  The cubic reciprocal gap from the direct
\(L^2\) current target is
\begin{equation}
 \label{eq:capacity-gap}
 \boxed{
 \frac{16}{21}-\frac12=\frac{11}{42}.
 }
\end{equation}
This is a no-go for the stated interpolation hull, not for structural
cancellation or Morrey estimates.

\begin{remark}[No pressure coercivity from cross density]
On \(\mathbb T^3\), choose \(k\neq0\), \(A\cdot k=0\), and
\[
 u(x)=a(x)=A\cos(k\cdot x).
\]
Then
\[
 (u\cdot\nabla)a=0,\qquad q=0,
 \qquad u\cdot a=|A|^2\cos^2(k\cdot x)\neq0.
\]
Thus the longitudinal pressure contraction cannot algebraically
coerce the trace of \(u\otimes a\).  This is a periodic symbol test,
not a same-trajectory Navier--Stokes counterexample.
\end{remark}

\section{Critical scalar comparison}

Qian and Xi prove an \(L^2\) energy inequality, uniqueness, a
fundamental solution, and regularity for scalar divergence-form
parabolic equations whose skew coefficient is uniformly bounded in
\(L^\infty_t\mathrm{BMO}_x\)~\cite{QXBMO}.  The q4 route does not
assume a uniform terminal coefficient bound, its datum is only a weak
trace, and its adjoint is a projected vector system.

Their separate mixed-Lebesgue paper treats a supercritical scalar
range but explicitly leaves supercritical regularity open
~\cite{QXMixed}.  Neither theorem supplies
Corollary~\ref{cor:gain}.

\section{What changed in this round}

\subsection*{Formal findings, subject to external review}

\begin{enumerate}
\item The primal--adjoint cross density has the exact local current
\eqref{eq:current} and pressure functional
\eqref{eq:pressure-payment}.
\item Its nonzero terminal measure is outside \(H^{-1}\), forcing
\eqref{eq:dyadic-divergence}--\eqref{eq:preterminal-divergence}.
\item Fourier localisation cancels both pressures and forces the signed
commutator cascade \eqref{eq:spectral-cascade}.
\item Finite weak-\(L^3\)-weighted adjoint dissipation renormalises the
reverse equation; hence the survivor forces
\eqref{eq:weighted-divergence}.
\item The temporal Lorentz gain \eqref{eq:gain} closes q4.
\item The optimised cubic current approaches, but does not claim, the
\(L^{21/16,7/5}\) endpoint and misses the direct \(L^2\)-capacity
target by \eqref{eq:capacity-gap}.
\item Cross density alone cannot coerce the adjoint pressure.
\end{enumerate}

\subsection*{Closed shortcuts}

\begin{enumerate}
\item A pressure lower bound cannot be inferred from the cross density's
nonzero trace.
\item Direct Hölder--Sobolev current interpolation does not reach the
\(\HH^1\)-null capacity.
\item The entrance cannot have finite
\(\int M\norm{\nabla a}_2^2\); any proposed strong-trace proof that
implicitly assumes this must prove it.
\item Global spectral localisation cannot expose a pressure payment,
because pressure cancels exactly there.
\end{enumerate}

\subsection*{Open obligations}

\begin{enumerate}
\item Prove \(\nabla a\in L_t^{22/9,2}L_x^2\), or directly prove
finite weighted adjoint dissipation from self-generation of the drift.
\item Alternatively prove equiintegrability of the commutators in
\eqref{eq:spectral-cascade}.
\item Find a cancellation-sensitive Morrey or total-current estimate
that reaches the \(\HH^1\)-null capacity.
\item Extend any closure beyond exact q4 to slower clocks and divergent
normalised energy.
\item Prove one complete Clay alternative for arbitrary admissible data.
\end{enumerate}

\begin{conjecture}[Self-generated reverse H\"older gain]
Under the exact q4 hypotheses, the same-trajectory projected adjoint
satisfies
\[
 \nabla a\in L_t^{22/9,2}L_x^2
\]
on some terminal interval.
\end{conjecture}

Corollary~\ref{cor:gain} shows that this conjecture excludes the q4
terminal energy defect.  It is not proved.

\section*{Source provenance}

The pinned arXiv v2 archive for the critical BMO comparison has
SHA-256
\[
\texttt{78af05bac214145c5b6d54c297867ab5f7d7c80d60dc6c40c3ae5518cdb06df2}.
\]
Its source file
\texttt{Aronson-Nonsymmetric-2017-02-18.tex} has SHA-256
\[
\texttt{a8c65d4425496766a7acb2d25b30603f4e79fadf5f400d92f2d65da206538a36}.
\]
The energy inequality and uniqueness theorem are at source lines
1220--1238, the continuity remark at lines 1332--1337, and the
initial-value theorem at lines 1341--1372.

\begin{thebibliography}{9}
\raggedright

\bibitem{QXBMO}
Z.~Qian and G.~Xi,
\emph{Parabolic equations with singular divergence-free drift vector
fields},
Journal of the London Mathematical Society 100 (2019), 17--40,
\href{https://doi.org/10.1112/jlms.12202}
{doi:10.1112/jlms.12202}.

\bibitem{QXMixed}
Z.~Qian and G.~Xi,
\emph{Parabolic equations with divergence-free drift in space
\(L_t^\ell L_x^q\)},
Indiana University Mathematics Journal 68 (2019), 761--797,
\href{https://doi.org/10.1512/iumj.2019.68.7685}
{doi:10.1512/iumj.2019.68.7685}.

\end{thebibliography}

\end{document}
